
La
démonstration se fait par récurrence sur la dimension @@M0@@ du support de @@M1@@; lorsque @@M2@@, c'est trivial. Prenons une section @@M3@@ qui ne s'annule en aucun point de l'ensemble fini @@M4@@; alors @@M5@@ définit une suite exacte @@M6@@, et @@M7@@ est @@M8@@-régulier par 1.2., à support de dimension @@M9@@.

Considérons la suite exacte correspondante :
@@M10@@
Par l'hypothèse de récurrence, on a @@M11@@ pour @@M12@@, donc @@M13@@ implique @@M14@@; mais @@M15@@, ce qui donne (i).

Pour établir (ii), on fait la chasse dans le diagramme commutatif suivant :
@@M16@@
où la flèche oblique est définie par @@M17@@, les lignes horizontales sont exactes, la première par la @@M18@@-régularité de @@M19@@ et par (i), déjà prouvé, et la seconde colonne est exacte par l'hypothèse de récurrence.

Enfin, (iii) découle du diagramme commutatif suivant :
@@M20@@
